\chapter{Results and Discussion}
\label{chap:results}

This chapter presents and interprets the 
empirical results across all experimental 
settings. Section~\ref{sec:primary} reports 
the primary result on the German Credit 
low-data regime, directly testing H1 and H3. 
Section~\ref{sec:null_control} presents the 
null control on the full-data regime, testing 
H2. Section~\ref{sec:radon} presents the 
Minnesota Radon supplementary null control. 
Section~\ref{sec:ablation_shared} presents 
the architectural ablation validating the 
shared backbone. Section~\ref{sec:supplementary} 
presents the diagnostic analyses: prior 
sensitivity, weight stability, 
leave-one-prior-out, calibration, and 
optimiser convergence. 
Section~\ref{sec:synthetic} reports the 
synthetic 2D results. 
Section~\ref{sec:hypothesis_summary} 
summarises the hypothesis outcomes.
Section~\ref{sec:fm_vs_cnf} presents a 
direct speed comparison between flow 
matching and CNF training, validating 
the simulation-free computational advantage.

\section{Primary Result: German Credit 
Low-Data Regime}
\label{sec:primary}

The primary experimental setting is German 
Credit with $n = 75$ training observations 
and 25 parameters. Under this 3:1 
data-to-parameter ratio the likelihood 
is insufficiently informative to overwhelm 
any reasonable prior, placing the experiment 
firmly in the prior-sensitive regime as 
defined in Section~3.1: the gap between 
the best and worst prior on the validation 
set reaches 0.13 to 0.27 nats across seeds, 
well above the 0.01 nat threshold.

\subsection{Mean NLL and Regret}

Table~\ref{tab:german_low_nll} reports the 
test NLL per seed and the mean across 5 
seeds for all eight weight optimisation 
algorithms, together with the oracle 
best-prior NLL computed post-hoc on the 
test set. Table~\ref{tab:german_low_regret} 
reports the regret relative to the oracle 
at each seed, with the maximum and mean 
regret across seeds.

\begin{table}[h]
\centering
\caption{Test NLL per seed and mean for all 
algorithms on German Credit low-data regime 
($n = 75$, 5 seeds). Lower is better.}
\label{tab:german_low_nll}
{\singlespacing
\begin{tabular}{lrrrrrr}
\hline
Algorithm & Seed 0 & Seed 1 & Seed 2 
& Seed 3 & Seed 4 & Mean \\
\hline
EM          & 0.6307 & 0.6214 & 0.6169 
             & 0.5976 & 0.6461 & \textbf{0.6225} \\
Proj. GD    & 0.6310 & 0.6215 & 0.6176 
             & 0.5972 & 0.6464 & 0.6227 \\
Coord. pair & 0.6312 & 0.6215 & 0.6175 
             & 0.5971 & 0.6468 & 0.6228 \\
Grid search & 0.6312 & 0.6215 & 0.6175 
             & 0.5971 & 0.6467 & 0.6228 \\
Frank-Wolfe & 0.6312 & 0.6215 & 0.6175 
             & 0.5971 & 0.6468 & 0.6228 \\
Best-single & 0.6333 & 0.6222 & 0.6175 
             & 0.5971 & 0.6500 & 0.6240 \\
Mirror desc.& 0.6380 & 0.6324 & 0.6258 
             & 0.6115 & 0.6546 & 0.6325 \\
Uniform     & 0.6753 & 0.6735 & 0.6671 
             & 0.6473 & 0.6881 & 0.6703 \\
\hline
Oracle      & 0.6307 & 0.6214 & 0.6169 
             & 0.5971 & 0.6461 & 0.6224 \\
\hline
\end{tabular}}
\end{table}

\begin{table}[h]
\centering
\caption{Regret (NLL minus oracle NLL) per 
seed for all algorithms. Max and mean regret 
across 5 seeds reported. Lower is better.}
\label{tab:german_low_regret}
{\singlespacing
\begin{tabular}{lrrrrrrrr}
\hline
Algorithm & Seed 0 & Seed 1 & Seed 2 
& Seed 3 & Seed 4 & Max & Mean \\
\hline
EM          & 0.0000 & 0.0000 & 0.0000 
             & 0.0005 & 0.0000 & \textbf{0.0005} & 0.0001 \\
Grid        & 0.0005 & 0.0001 & 0.0006 
             & 0.0000 & 0.0006 & 0.0006 & 0.0004 \\
Proj. GD    & 0.0003 & 0.0001 & 0.0007 
             & 0.0001 & 0.0003 & 0.0007 & 0.0003 \\
Coord. pair & 0.0005 & 0.0001 & 0.0006 
             & 0.0000 & 0.0007 & 0.0007 & 0.0004 \\
Frank-Wolfe & 0.0005 & 0.0001 & 0.0006 
             & 0.0000 & 0.0007 & 0.0007 & 0.0004 \\
Best-single & 0.0026 & 0.0008 & 0.0006 
             & 0.0000 & 0.0039 & 0.0039 & 0.0016 \\
Mirror desc.& 0.0073 & 0.0110 & 0.0089 
             & 0.0144 & 0.0085 & 0.0144 & 0.0100 \\
Uniform     & 0.0446 & 0.0521 & 0.0502 
             & 0.0502 & 0.0420 & 0.0521 & 0.0478 \\
\hline
\end{tabular}}
\end{table}

\subsection{H1: Guarantee Visibility}
\label{sec:h1}

All principled mixture algorithms (EM, 
Proj. GD, Coord., Grid, Frank-Wolfe) 
achieve lower mean test NLL than 
best-single: 0.6225--0.6228 versus 
0.6240. More importantly, the worst-case 
regret of the mixture algorithms is 
substantially lower. EM achieves a maximum 
regret of 0.0005 nats, compared to 0.0039 
nats for best-single, a factor of 
approximately eight. The source of 
best-single's regret is clear from 
Table~\ref{tab:prior_per_seed}: on seeds 
0, 1, and 2 the best prior is Gaussian, 
while on seeds 3 and 4 it is Laplace. A 
practitioner who commits to Gaussian in 
advance incurs regret up to 0.003 nats; 
one who commits to Laplace incurs regret 
up to 0.012 nats. The credal mixture 
avoids this entirely by distributing 
weight across the credal set and 
adapting to whichever prior is most 
informative for each data subsample.

\begin{table}[h]
\centering
\caption{Per-prior test predictive 
log-likelihood per seed (German Credit, 
$n = 75$). Bold indicates the best prior 
for each seed.}
\label{tab:prior_per_seed}
{\singlespacing
\begin{tabular}{lrrrrr}
\hline
Prior & Seed 0 & Seed 1 & Seed 2 
& Seed 3 & Seed 4 \\
\hline
Gaussian      & \textbf{$-$0.630} 
              & \textbf{$-$0.622} 
              & \textbf{$-$0.616} 
              & $-$0.609 & $-$0.650 \\
Laplace       & $-$0.633 & $-$0.624 
              & $-$0.618 
              & \textbf{$-$0.597} 
              & \textbf{$-$0.644} \\
Cauchy        & $-$0.695 & $-$0.690 
              & $-$0.711 & $-$0.685 
              & $-$0.725 \\
Student-T     & $-$0.718 & $-$0.711 
              & $-$0.696 & $-$0.683 
              & $-$0.716 \\
GaussianWide  & $-$0.823 & $-$0.888 
              & $-$0.853 & $-$0.793 
              & $-$0.790 \\
\hline
\end{tabular}}
\end{table}

This instability in which prior is 
best across seeds is the core motivation 
for the credal mixture approach, and 
Table~\ref{tab:prior_per_seed} makes 
it visible concretely. H1 is 
\textbf{confirmed}: in the 
prior-sensitive regime, the credal 
mixture achieves lower worst-case 
regret than the best-single 
strategy.

\subsection{H3: Algorithm Comparison}
\label{sec:h3}

EM achieves both the lowest mean NLL 
(0.6225) and the lowest maximum regret 
(0.0005 nats) among all algorithms. 
The other principled optimisers (Proj. 
GD, Coord., Grid, Frank-Wolfe) are 
essentially tied on mean NLL (0.6227--0.6228) 
and close on max regret (0.0006--0.0007). 
EM's advantage is most visible at 
seed 3, where it achieves near-oracle 
performance (regret 0.0005) while 
Grid and Coord. achieve exactly oracle 
(regret 0.0000), this is due to 
different initialisation behaviour 
on that particular subsample.

Mirror descent performs considerably 
worse (mean NLL 0.6325, max regret 
0.0144), consistent with its known 
sensitivity to the learning rate 
($\alpha = 0.05$). Uniform weighting 
is worst by a large margin (mean 
0.6703), reflecting that priors 
vary substantially in quality and 
equal weighting dilutes the 
contribution of the best components. 
H3 is \textbf{confirmed}: EM achieves 
the lowest worst-case regret. Its 
advantage over alternatives is further 
supported by its closed-form M-step, 
which requires no learning rate 
tuning and converges reliably within 
200 iterations on every seed.

\subsection{MCMC Reference Baseline}

To confirm that the prior-sensitivity 
pattern is real rather than an 
artefact of the flow approximation, 
the same evaluation is performed 
using HMC posterior samples in place 
of flow samples. Under exact HMC 
posteriors, Gaussian achieves 
predictive log-likelihood of 
approximately $-0.616$ to $-0.630$ 
nats across seeds while GaussianWide 
achieves $-0.790$ to $-0.888$ nats, 
a gap of 0.16 to 0.26 nats that 
closely mirrors the pattern in 
Table~\ref{tab:prior_per_seed}. The 
flow approximation adds a consistent 
overhead of approximately 0.015 nats 
relative to HMC across all priors, 
confirming that the approximation 
error is small and approximately 
uniform across the credal set. This 
is the empirical basis for 
Assumption~\ref{ass:transfer_ch2}: 
the transfer error $\eta$ is small 
and consistent, not concentrated 
on particular components. Connecting 
this to Proposition~\ref{prop:sqrt_eps_ch2}, 
the uniform 0.015 nat overhead 
provides an empirical upper bound 
on $\eta$, suggesting the bounded 
density ratio condition ($\Gamma$) 
and log-ratio condition ($M$) of 
the proposition are approximately 
satisfied in this setting, even 
though they are not verified 
analytically.

\section{Null Control: German Credit 
Full-Data Regime}
\label{sec:null_control}

Table~\ref{tab:german_full} reports 
results for the full-data setting 
($n = 700$, 28 observations per 
parameter). At this data-to-parameter 
ratio the likelihood dominates and 
all posteriors converge regardless 
of prior.

\begin{table}[h]
\centering
\caption{Test NLL for all algorithms 
on German Credit full-data regime 
($n = 700$, 10 runs). All algorithms 
are within 0.002 nats.}
\label{tab:german_full}
{\singlespacing
\begin{tabular}{lrr}
\hline
Algorithm & Mean NLL & Std \\
\hline
Proj. GD    & \textbf{0.4981} & 0.0003 \\
Grid search & 0.4985 & 0.0007 \\
Coord. pair & 0.4985 & 0.0007 \\
Frank-Wolfe & 0.4985 & 0.0007 \\
Best-single & 0.4990 & 0.0011 \\
EM          & 0.4991 & 0.0001 \\
Mirror desc.& 0.4999 & 0.0001 \\
Uniform     & 0.5000 & 0.0001 \\
\hline
\end{tabular}}
\end{table}

The total spread across all algorithms 
is 0.002 nats, well within the 0.005 
nat threshold defined for H2. Best-single 
converges to a single dominant prior 
on all 10 runs, confirming that the 
data is sufficient to identify the 
posterior regardless of prior. The 
Wilcoxon signed-rank test on the 
paired difference in test NLL between 
EM and best-single across runs yields 
$p > 0.05$, confirming no statistically 
significant difference. H2 is 
\textbf{confirmed}: the credal mixture 
provides no benefit and no harm when 
the data is sufficient to overwhelm 
prior influence. This validates the 
theoretical prediction from 
Corollary~\ref{cor:recovery_ch2} 
when all posteriors have converged, 
the credal mixture reduces to the 
best single component.

\section{Supplementary Null Control: 
Minnesota Radon}
\label{sec:radon}

Table~\ref{tab:radon} reports results 
for the Radon experiment (2D, 
$K = 3$ priors, 10 seeds). With a 
310:1 data-to-parameter ratio, this 
dataset is deeply in the 
prior-insensitive regime.

\begin{table}[h]
\centering
\caption{Test NLL for all algorithms 
on Minnesota Radon (10 seeds). 
All principled algorithms are 
within 0.001 nats of best-single.}
\label{tab:radon}
{\singlespacing
\begin{tabular}{lrr}
\hline
Algorithm & Mean NLL & Std \\
\hline
Best-single & \textbf{2.5141} & 0.0085 \\
Grid search & 2.5141 & 0.0085 \\
Coord. pair & 2.5141 & 0.0085 \\
Proj. GD    & 2.5141 & 0.0084 \\
Frank-Wolfe & 2.5141 & 0.0085 \\
EM          & 2.5146 & 0.0085 \\
Mirror desc.& 2.5224 & 0.0087 \\
Uniform     & 2.5257 & 0.0088 \\
\hline
\end{tabular}}
\end{table}

All principled algorithms (EM, Grid, 
Coord., Proj. GD, Frank-Wolfe) achieve 
essentially identical NLL to best-single, 
confirming the null result. The Gaussian 
prior dominates completely: grid search 
and coord. assign weight $\approx 1.0$ 
to Gaussian, and EM gives it 88\% weight. 
This is consistent with a linear 
regression model under abundant data: 
the 310:1 ratio means the prior has 
negligible influence and the mixture 
provides no additional information. 
Mirror descent and uniform are 
considerably worse (2.5224 and 2.5257 
respectively), reflecting their failure 
to concentrate weight on the dominant 
component.

\section{Architectural Ablation: Shared 
vs Separate Models}
\label{sec:ablation_shared}

Figure~\ref{fig:ablation} compares the 
shared-backbone architecture against five 
separately trained flows, one per prior, 
under uniform mixture weights. This 
isolates the effect of the architecture 
from the weight selection procedure.

\begin{figure}[h]
\centering
\includegraphics[width=\textwidth]
{exp6_ablation_shared_vs_separate}
\caption{Architectural ablation: shared 
conditional model vs separately trained 
flows. Left: training loss curves. 
Centre: predictive log-likelihood under 
uniform mixture weights ($-0.4929$ vs 
$-0.5028$). Right: parameter count 
(888.6k vs 2665.7k, a 67\% reduction 
with the shared model).}
\label{fig:ablation}
\end{figure}

The shared model achieves predictive 
log-likelihood of $-0.4929$ nats under 
uniform weights, compared to $-0.5028$ 
nats for the separate models, an 
improvement of 0.0099 nats. The shared 
model also converges faster and with 
lower training loss variance, visible 
in the left panel. The parameter count 
is reduced by 67\% (888.6k vs 2665.7k 
parameters), as the shared backbone 
amortises the cost of representing the 
velocity field across all $K$ priors 
simultaneously rather than maintaining 
$K$ independent sets of parameters.

The result directly validates 
Assumption~\ref{ass:transfer_ch2}. 
The separate training approach cannot 
guarantee comparable approximation 
quality across priors, each model 
may converge to a different local 
optimum with different residual error 
$\varepsilon_k$. The shared backbone 
trains all $K$ posterior approximations 
jointly under the same optimisation 
trajectory, producing components with 
comparable approximation quality and 
a well-defined mixture. That it also 
achieves better predictive performance 
with fewer parameters confirms that 
the shared architecture is not merely 
a theoretical convenience but a 
practical improvement on all dimensions.

\section{Supplementary Analyses}
\label{sec:supplementary}

\subsection{Prior Sensitivity}
\label{sec:prior_sensitivity}

Figure~\ref{fig:prior_sensitivity} 
quantifies the degree to which each 
prior shifts the posterior mean 
relative to the Gaussian baseline, 
across all 25 parameter dimensions.

\begin{figure}[h]
\centering
\includegraphics[width=\textwidth]
{exp1_prior_sensitivity_curve}
\caption{Prior sensitivity analysis on 
German Credit ($n = 75$). Left: 
posterior mean shift per dimension 
relative to the Gaussian prior. 
Right: maximum absolute mean shift 
across dimensions per prior. The 
credal mixture (dashed line, 0.004) 
is an order of magnitude more stable 
than the most sensitive individual 
prior (GaussianWide, 0.028).}
\label{fig:prior_sensitivity}
\end{figure}

GaussianWide produces the largest 
posterior shift (max 0.028), followed 
by Laplace (0.027), Cauchy (0.021), 
and Student-T (0.020). Gaussian, as 
the reference, has zero shift by 
definition. The credal mixture achieves 
a maximum mean shift of only 0.004 
across all dimensions, substantially 
lower than any individual component. 
This confirms that the mixture posterior 
is more stable than any constituent 
prior, which is the operational meaning 
of prior robustness: rather than 
committing to the posterior shifts 
induced by a particular prior, the 
credal mixture averages out 
prior-induced bias.

The prior sensitivity figure also 
provides visual confirmation that the 
prior-sensitive regime is genuine. 
The individual prior curves in the 
left panel show substantial 
dimension-by-dimension disagreement, 
particularly on dimensions 5 through 
15 where the Laplace and GaussianWide 
curves diverge significantly from 
Gaussian. This is the geometric 
manifestation of the NLL gap reported 
in Table~\ref{tab:prior_per_seed}.

\subsection{Weight Stability}
\label{sec:weight_stability}

Figure~\ref{fig:weight_stability} 
shows the distribution of EM mixture 
weights across 10 seeds on the German 
Credit low-data regime.

\begin{figure}[h]
\centering
\includegraphics[width=\textwidth]
{exp2_weight_stability}
\caption{EM mixture weight distribution 
across 10 seeds. Left: box plots showing 
mean and variance per prior (Cauchy: 
$0.26 \pm 0.194$; Gaussian: 
$0.24 \pm 0.164$; Laplace: 
$0.19 \pm 0.165$; Student-T: 
$0.16 \pm 0.126$; GaussianWide: 
$0.15 \pm 0.126$). Right: per-seed 
weight heatmap showing that the 
dominant prior changes across 
data subsamples.}
\label{fig:weight_stability}
\end{figure}

The large standard deviations 
(0.126--0.194) confirm that weight 
assignment is genuinely unstable 
across seeds: no single prior 
consistently dominates, and the 
credal mixture must adapt its 
weighting to each data subsample. 
Seed 1 is a notable outlier where 
Cauchy receives weight approximately 
0.65, a heavy-tailed subsample where 
the Cauchy prior is substantially 
better calibrated. The maximum weight 
standard deviation across components 
is 0.194, the highest value across all 
experiments, confirming that the 
low-data regime is genuinely 
prior-sensitive and that the instability 
is not a numerical artefact. The 
adaptive reweighting is precisely 
the mechanism by which the credal 
mixture achieves near-oracle performance 
across seeds despite the identity of 
the best prior changing between seeds.

\subsection{Leave-One-Prior-Out}
\label{sec:lopo}

The Leave-One-Prior-Out (LOPO) 
analysis removes each prior in turn 
from the full five-component credal 
set and evaluates the resulting 
four-component mixture. 
Figure~\ref{fig:lopo} shows the 
change in test predictive 
log-likelihood relative to the full 
five-component mixture.

\begin{figure}[h]
\centering
\includegraphics[width=\textwidth]
{exp3_lopo}
\caption{Leave-One-Prior-Out analysis. 
Left: change in predictive 
log-likelihood when each prior is 
removed ($\Delta > 0$ means the 
four-component mixture is better). 
Right: absolute predictive 
log-likelihood for all configurations.}
\label{fig:lopo}
\end{figure}

Removing GaussianWide or Gaussian 
improves the mixture by $+0.00040$ 
and $+0.00037$ nats respectively, 
confirming that these two components 
are net detractors when averaged 
across seeds. Removing Laplace 
($-0.00036$), Student-T ($-0.00038$), 
or Cauchy ($-0.00038$) degrades 
performance, confirming these are 
genuinely useful contributors. All 
magnitudes are below 0.001 nats, 
indicating that no single prior is 
critical to the mixture's robustness 
and the credal set is not far from 
the optimal subset.

The marginal degradation from including 
GaussianWide and Gaussian is a 
predictable consequence of their 
behaviour in the prior-sensitive regime. 
GaussianWide is a weakly informative 
prior that produces a diffuse posterior, 
it is included in the credal set 
to represent the range of plausible 
prior beliefs, but it rarely produces 
the best posterior and its inclusion 
dilutes the mixture. This is not a 
failure of the method: the theoretical 
guarantee of 
Theorem~\ref{thm:exact_mixture_ch2} 
says the mixture will be at least as 
good as the best component in 
expectation, and including extra 
poor-quality components introduces a 
small cost that the weight optimisation 
does not fully correct. This finding 
raises a practical tension: a 
practitioner could improve predictive 
performance by pruning GaussianWide 
from the credal set, but doing so 
requires knowing in advance which 
priors are net detractors, which 
is precisely the prior knowledge 
the credal set is designed to avoid 
assuming. Automatic credal set 
selection remains an open problem.

\subsection{Calibration}
\label{sec:calibration}

Figure~\ref{fig:calibration} shows 
calibration curves and Expected 
Calibration Error (ECE) for each 
method on the German Credit full-data 
setting.

\begin{figure}[h]
\centering
\includegraphics[width=\textwidth]
{exp4_calibration}
\caption{Calibration analysis on German 
Credit full-data regime. Left: 
reliability diagrams. Right: ECE 
and Brier score per method.}
\label{fig:calibration}
\end{figure}

ECE values range from 0.047 (Student-T, 
best) to 0.069 (Laplace, worst). The 
credal mixture achieves ECE of 0.052, 
comparable to Gaussian (0.053) and 
better than GaussianWide (0.059) and 
Laplace (0.069). Brier scores are 
essentially identical across all 
methods (0.165--0.166), confirming 
that the methods differ in systematic 
bias (captured by ECE) rather than 
in overall sharpness. The calibration 
curves show that all methods track the 
diagonal closely in the moderate 
probability region (0.2--0.8), with 
the largest deviations at the extremes 
where test counts are low. 

The credal mixture does not improve 
on the best individual prior for 
calibration in the full-data regime, 
which is consistent with H2: when 
all posteriors have converged, 
combination provides no additional 
benefit. The result confirms that 
the mixture does not degrade 
calibration either, which is a 
practically relevant finding, a 
practitioner adopting the credal 
mixture approach incurs no calibration 
penalty in the data-rich regime.

\subsection{Optimiser Convergence}
\label{sec:optimiser}

Figure~\ref{fig:optimiser} compares 
convergence trajectories and final 
weight vectors for PGD, EM, and 
Frank-Wolfe on the full-data German 
Credit setting.

\begin{figure}[h]
\centering
\includegraphics[width=\textwidth]
{exp5_optimiser_comparison}
\caption{Optimiser convergence on 
German Credit full-data regime. Left: 
convergence of mixture predictive 
log-likelihood across iterations. 
Right: final weight vectors for each 
algorithm. PGD and EM converge to 
near-identical solutions 
($\|\Delta w\|_2 < 0.01$) while 
Frank-Wolfe degenerates to a sparse 
solution assigning weight $\approx 1.0$ 
to Student-T.}
\label{fig:optimiser}
\end{figure}

PGD and EM converge to nearly 
identical objectives and weight 
vectors: the maximum pairwise 
$\|\Delta w\|_2$ between their 
final solutions is 0.119, with both 
placing substantial weight on 
GaussianWide (0.426 and 0.368), 
Gaussian (0.278 and 0.251), and 
Student-T (0.268 and 0.368). This 
convergence to a similar interior 
point of the simplex confirms that 
the NLL objective is well-conditioned 
and the solution is stable across 
optimisation algorithms.

Frank-Wolfe behaves qualitatively 
differently. Its harmonic step size 
schedule drives it toward a sparse 
vertex solution: it assigns weight 
1.0 to Student-T and 0.0 to all 
other priors, effectively degenerating 
to single-component selection. This 
explains why Frank-Wolfe achieves 
the same NLL as best-single in the 
full-data regime (Table~\ref{tab:german_full}) 
and why it does not exploit the 
mixture structure. In the prior-sensitive 
low-data regime where best-single 
already achieves low regret on seed 3, 
Frank-Wolfe also achieves low regret 
(0.0007 max, matching the other 
algorithms), but this is because 
sparse selection happens to select 
the right component, not because 
the mixture is genuinely being used. 
Practitioners who require an 
interior-point credal mixture should 
avoid Frank-Wolfe and prefer EM or 
Proj. GD, both of which consistently 
produce non-degenerate weight vectors.

\subsection{Flow Matching vs CNF Training Speed}
\label{sec:fm_vs_cnf}

Figure~\ref{fig:speed_comparison} compares 
the per-step wall-clock training time of 
flow matching against a CNF baseline using 
an identical \texttt{VelocityResNet} 
architecture on German Credit (seed 0).
The CNF baseline measures the cost of 
10-step Euler integration of the augmented 
ODE with Hutchinson trace estimation, 
which is the minimum simulation cost 
required for exact CNF density evaluation 
during training.

\begin{figure}[h]
\centering
\includegraphics[width=\textwidth]
{plot_speed_comparison}
\caption{Flow matching vs CNF training 
speed comparison. Left: per-step time 
distributions over 200 trials. Centre: 
flow matching training loss curve (CNF 
not trained to convergence). Right: 
test NLL per prior under flow matching.}
\label{fig:speed_comparison}
\end{figure}

Table~\ref{tab:speed_comparison} 
summarises the results.

\begin{table}[h]
\centering
\caption{Speed comparison between flow 
matching and CNF baseline on German Credit 
low-data regime ($n = 75$, seed 0, 
6{,}000 steps, identical 
\texttt{VelocityResNet} architecture). 
CNF total time is estimated from per-step 
timing; the CNF was not trained to 
convergence.}
\label{tab:speed_comparison}
{\singlespacing
\begin{tabular}{lrr}
\hline
Metric & Flow Matching & CNF Baseline \\
\hline
Total training time (s) & 106.7 & 1{,}277.0 \\
Mean step time (ms) & 11.7 & 139.7 \\
Speedup & \multicolumn{2}{c}{$12.0\times$} \\
\hline
\end{tabular}}
\end{table}

Flow matching is 12x faster per training 
step. The CNF baseline requires 10 
velocity field evaluations per training 
step for ODE integration, plus Hutchinson 
trace estimation at each step, while flow 
matching requires a single forward pass 
and a regression loss. Projected over 
6{,}000 training steps, CNF training 
would require approximately 21 minutes 
versus 107 seconds for flow matching on 
the same hardware. In the credal mixture 
setting where $K = 5$ posteriors are 
approximated, this gap compounds: 
independent CNF training across all 
five priors would require approximately 
105 minutes versus the 107 seconds of 
the shared flow matching model. This 
directly validates the computational 
motivation of Section~2.5.3.

\section{Synthetic 2D Datasets}
\label{sec:synthetic}

The three synthetic datasets, 
Eight-Ring, Two Moons, and Two-Arm 
Spirals, test whether the shared 
conditional flow can represent complex 
multimodal density structure. These 
experiments do not test H1, H2, or 
H3 directly; rather, they verify the 
geometric prerequisite that the flow 
architecture is expressive enough to 
approximate complex posteriors in 
the first place.

Table~\ref{tab:synthetic} reports 
mean test KDE-NLL across 10 seeds 
for all algorithms, and the gap 
between best-single and the 
coordinate descent mixture on each 
dataset.

\begin{table}[h]
\centering
\caption{Mean test KDE-NLL (lower is 
better) across 10 seeds for synthetic 
2D datasets. Gap is best-single minus 
credal (coord.) KDE-NLL; positive 
means credal is better.}
\label{tab:synthetic}
{\singlespacing
\begin{tabular}{lrrrr}
\hline
Dataset & Uniform & Best-single 
& Credal (coord.) & Gap \\
\hline
Eight-Ring   & 2.018 & 1.983 
             & 1.982 & 0.001 \\
Two Moons    & 2.019 & 1.982 
             & 1.981 & 0.001 \\
Two-Arm Spirals & 2.018 & 1.981 
             & 1.980 & 0.000 \\
\hline
\end{tabular}}
\end{table}

The credal mixture matches or 
marginally outperforms best-single 
on all three datasets (gaps of 
0.000--0.001 nats), consistent with 
the prior-insensitive regime: the 
synthetic datasets have effectively 
unlimited training data and the 
five prior perturbations produce 
similar approximations, so no 
combination benefit is expected. 
The uniform baseline is substantially 
worse (2.018--2.019 vs 1.980--1.983), 
confirming that the flow does learn 
informative density structure that 
a flat prior cannot capture.

The more substantive finding is that 
KDE-NLL values of approximately 
$1.98$ nats across all three datasets 
confirm that the shared conditional 
flow successfully approximates 
the target densities. The Eight-Ring 
geometry, with eight discrete modes 
on a circle, requires the flow to 
learn multi-modal structure with 
rotational symmetry. Two-Arm Spirals 
imposes a topologically complex 
density with long-range correlations. 
Two Moons requires capturing a 
non-linear bimodal structure. That 
the flow achieves comparable 
KDE-NLL across all three geometries 
demonstrates that the shared 
backbone has sufficient expressive 
capacity to serve as a reliable 
posterior approximator in the 
Bayesian experiments.

\section{Hypothesis Outcomes}
\label{sec:hypothesis_summary}

Table~\ref{tab:hypothesis_summary} 
summarises the outcome of each 
hypothesis against the key 
quantitative evidence.

\begin{table}[h]
\centering
\caption{Summary of hypothesis 
test outcomes.}
\label{tab:hypothesis_summary}
{\singlespacing
\begin{tabular}{p{3.5cm}lp{6.5cm}}
\hline
Hypothesis & Outcome & Key evidence \\
\hline
H1 (Guarantee Visibility) 
& \textbf{Confirmed} 
& EM max regret 0.0005 nats vs 
0.0039 for best-single; all 
principled mixture algorithms 
achieve lower mean NLL \\[4pt]
H2 (Prior Sensitivity Req.) 
& \textbf{Confirmed} 
& Full-data spread 0.002 nats 
across all algorithms; 
Wilcoxon $p > 0.05$ \\[4pt]
H3 (Algorithm Comparison) 
& \textbf{Confirmed} 
& EM achieves lowest max regret 
(0.0005) and lowest mean NLL 
(0.6225); closed-form, no 
learning rate \\
\hline
\end{tabular}}
\end{table}

All three hypotheses are confirmed. 
The KL convexity guarantee is 
empirically visible in the 
prior-sensitive low-data regime: 
the credal mixture achieves 
near-oracle worst-case regret of 
0.0005 nats under EM, compared 
to 0.0039 nats for best-single 
strategy, a factor of approximately 
eight improvement. The guarantee 
correctly disappears in the 
prior-insensitive regime, confirming 
that it is prior sensitivity, not 
the method itself, that drives the 
observed benefit. EM is the most 
consistent weight optimisation 
algorithm across the regret criterion 
and practical considerations. The 
shared-backbone architecture provides 
empirical validation for 
Assumption~\ref{ass:transfer_ch2}, 
outperforming separately trained 
flows on predictive performance, 
training stability, and parameter 
efficiency simultaneously.